\section{A smaller direct quartic with explicit coefficients}
\label{sec:direct-oriented}

The second completed construction in this report compiles a numerical proof
calculus directly.  Its final polynomial has small, completely materialized
integer coefficients.  The construction uses the calibrated $C_{222}$
calculus and a shared Chinese-remainder certificate; it does not require a
K215 transition table or Jones's universal polynomial.

Let $R_1,\ldots,R_{2089}$ be the literal quadratic residuals, in the
listed coordinate order, in \path{direct/eq/equation.json}. Define
\begin{equation}\label{eq:direct2473}
 P_{2473}(x_1,\ldots,x_{2473})
       =\sum_{j=1}^{2089}R_j(x_1,\ldots,x_{2473})^2.
\end{equation}
Every coefficient in every $R_j$ is a supplied integer. The companion
\path{direct/eq/expanded_polynomial.json} gives all $20\,580$ collected
monomials of the same polynomial, without any arithmetic macros.

\begin{theorem}[Direct consistency equation]
\label{thm:direct2473}
The polynomial $P_{2473}$ has exact degree four and coefficient height
$127\,808\,289<2^{27}$. Ordinary mathematics formalizable in ZFC proves
\[
 \bigl(\nexists x\in\mathbb N^{2473}\;P_{2473}(x)=0\bigr)
       \quad\longleftrightarrow\quad \operatorname{Con}(\mathrm{ZFC}),
 \qquad \mathbb N=\{0,1,2,\ldots\}.
\]
This theorem uses the identified C222/refutability calibration and the
Pell--Matiyasevich characterization. It does not assert either side of the
displayed equivalence.
\end{theorem}

The materialized incoming direct construction had $3\,614$ natural
coordinates and $2\,954$ quadratic residuals. The present construction
removes $1\,141$ coordinates, a $31.57\%$ reduction, and $865$ residuals.
Its expanded monomial count is larger. Thus it is a new point in a tradeoff
between coordinate count and expanded size, not a statement of global
minimality.

\begin{center}
\begin{tabular}{lrrrr}
\hline
Construction & Variables & Rows & Monomials & Height\\
\hline
Parent, affine gates removed & 3614 & 2954 & 15996 & 127778545\\
New, affine gates removed & 2473 & 2089 & 20580 & 127808289\\
Parent, affine gates retained & 6183 & 5523 & 17430 & 1993\\
New, affine gates retained & 4860 & 4476 & 14719 & 1837\\
\hline
\end{tabular}
\end{center}
All four rows describe exact quartics over nonnegative integers. The second
new equation, with retained affine gates, improves each displayed quantity
against its corresponding parent. Its exact data are under
\path{direct/small_height/}. The proofs of existence equivalence apply to
both variants.

\subsection{An eleven-coordinate graph for positive exponentiation}

The main new arithmetic fact is a compact positive-power certificate with
only eleven auxiliary natural coordinates before gate expansion. For
$a,x,y\in\mathbb N$, write
\[
 \mathcal N(a,x,y)=x^2+y^2-1-a^2y^2.
\]
The intended inputs are $n,k>0$ and the proposed output is $m\in\mathbb N$.
Take eleven auxiliary coordinates
\[
 U,\alpha,z,y,q,d,g,b,s,r,h\in\mathbb N.
\]
Define the following ordinary polynomial abbreviations:
\begin{align*}
 W&=n+k,& Z&=z+1,&Y&=k+y,\\
 A&=n+\alpha+1,&T&=m+g+1,&V&=Y^2(d+1),\\
 B&=1+4Y(q+1),&
 X&=Y(\alpha+1)+m+Th,\\
 S&=X+Us,&R&=k+4Yr.
\end{align*}
They are abbreviations, not extra quantified coordinates. The complete
graph consists of six equations:
\begin{align}
 \mathcal N(A,X,Y)&=0,&\mathcal N(A,U,V)&=0,
       &\mathcal N(B,S,R)&=0,\label{eq:op-norms}\\
 B&=A+Ub,&2An&=T+n^2+1,
       &\mathcal N(W+1,A,WZ)&=0.\label{eq:op-rest}
\end{align}

\begin{theorem}[Positive-power graph]
\label{thm:positive-power11}
For every $n,k>0$ and every $m\in\mathbb N$, equations
\eqref{eq:op-norms}--\eqref{eq:op-rest} have a solution in their eleven
auxiliary natural coordinates exactly when $m=n^k$.
\end{theorem}

\begin{proof}
We first recall the standard finite Pell--Matiyasevich criterion and then
prove each reduction to the displayed graph. The primary source is the
Mathlib theorem \texttt{Pell.eq\_pow\_of\_pell}, together with its
\texttt{matiyasevic} theorem for the Pell coordinate at a prescribed
index. The supplied source is pinned to Mathlib commit
\texttt{81a5d257c8e410db227a6665ed08f64fea08e997}; the repository's
\path{PellRepresentation.lean} and \path{PowerRepresentation.lean}
state the equivalent explicit finite arithmetic criteria.

For positive $n,k$, the original criterion existentially chooses
$W,A,T,Z,X,Y,U,V,S,R,B$ subject to
\begin{gather*}
 A\ge2,\quad Y\ge k,\quad B\ge2,\quad
 \mathcal N(A,X,Y)=\mathcal N(A,U,V)=\mathcal N(B,S,R)=0,\\
 B\equiv1\pmod{4Y},\qquad V\in Y^2\mathbb Z_{>0},\\
 B\equiv A\pmod U,\quad S\equiv X\pmod U,
       \quad R\equiv k\pmod{4Y},\\
 X+Yn\equiv YA+m\pmod T,\quad
 2An=T+n^2+1,\quad T>m,\\
 W\ge n,k,\qquad \mathcal N(W+1,A,WZ)=0.
\end{gather*}
These conditions are equivalent to $m=n^k$. In particular, the modular
conditions are not informal substitutes for a finite graph: they will
be replaced by the explicit natural quotients below.

First we justify orienting all four congruences. Since
$V\ge Y^2\ge Y\ge1$, the first two norms imply
\[
 U\ge X,\qquad U\ge A.
\]
Also $B>0$ and $S>0$, the latter following from its norm. If
$0<a\le u$, every positive representative of $a$ modulo $u$ is
$a+u j$ for some $j\in\mathbb N$. This includes the boundary $a=u$,
where the least positive representative is $u$. Hence
\[
 B=A+Ub,\qquad S=X+Us
\]
for natural $b,s$. As $0<k\le Y<4Y$, one similarly has
$R=k+4Yr$ for a natural $r$.

For the last congruence, the first norm gives
\[
 X^2-(A-1)^2Y^2=1+2(A-1)Y^2>0,
\]
and thus $X>(A-1)Y$. Therefore
\[
 D:=X+Yn-YA-m>(n-1)Y-m\ge-m>-T.
\]
Because $D$ is an integer multiple of positive $T$, it cannot be
negative. Consequently
\begin{equation}\label{eq:power-natural-quotient}
 X+Yn=YA+m+Th\qquad(h\in\mathbb N).
\end{equation}
This interval argument, rather than an assumption about the signs of
formal quotient witnesses, licenses the fourth orientation.

Next we fix $W=n+k$. The standard completeness proof chooses
$W=\max(n,k)$ but uses only $W\ge n,k$ and $W>0$ after that choice.
Here is the actual construction for any such $W$. For the Pell sequence
\[
 x_j(a)+y_j(a)\sqrt{a^2-1}
       =(a+\sqrt{a^2-1})^j,
\]
put $A=x_W(W+1)$ and $Z=y_W(W+1)/W$. The recurrence congruence
$y_j(a)\equiv j\pmod{a-1}$ proves that $Z$ is an integer. It is
positive. Its Pell equation is exactly the last norm, and the usual
elementary growth estimate gives
\[
 n^k\le n^W<(W+1)^W\le A.
\]
Set $X=x_k(A)$, $Y=y_k(A)$, and $T=2An-n^2-1$. Then $Y\ge k$.
Moreover $A\ge n+1$ gives
\[
 T-A=(2n-1)A-n^2-1\ge n^2+n-2\ge0,
\]
so $T\ge A>n^k$. The identity
$A^2-1\equiv(n-A)^2\pmod T$, evaluated in the quadratic-ring
recurrence, gives $X+Yn\equiv YA+n^k\pmod T$. The standard
Pell-index theorem supplies the remaining $U,V,S,R,B$ for this fixed
$A,X,Y,k$. Thus $W=n+k$ preserves completeness.

The original graph with $W=n+k$ has $A\ge2$ and
\[
 A^2=1+W^3(W+2)Z^2.
\]
Hence $Z\ge1$ and $A>W\ge n$. We may write
\[
 Z=z+1,\quad A=n+\alpha+1,\quad Y=k+y,
 \quad T=m+g+1,\quad V=Y^2(d+1).
\]
Since $B\ge2$ and $B\equiv1\pmod{4Y}$, the positive quotient
$(B-1)/(4Y)$ is at least one; write $B=1+4Y(q+1)$.
All the new coordinates are natural. Equation
\eqref{eq:power-natural-quotient} now solves for
$X=Y(\alpha+1)+m+Th$, and the other two oriented equations solve
for $S=X+Us$ and $R=k+4Yr$. Substituting these three nonnegative
polynomial expressions removes three coordinates and three equations.
The result is precisely the eleven-coordinate graph in the theorem.

Conversely, from a solution of that graph, $W\ge2$ and $Z\ge1$
force $A>W\ge n$ through the last norm. All original inequalities,
norms and congruences are then restored by the displayed definitions.
The original Pell--Matiyasevich criterion gives $m=n^k$.
\end{proof}

\subsection{Why positive inputs suffice at every use}

The completed direct compiler has exactly $55$ power graphs, $16$
binomial graphs and two factorial graphs. Its outer length and label
bound are parameterized by $t=t_0+1$ and $m_0+1$, respectively.
Every encoded CRT coordinate is also an independent natural coordinate
plus one. These choices preserve completeness because the Chinese
remainder theorem specifies residue classes, and representatives can
be increased by positive multiples of their common modulus.

All bases and exponents of the actual power calls are positive. The
carrier has positive base and exponent $t>0$. Each range-binomial
graph has positive upper argument and positive lower argument. The
factorial arguments are positive and their radius arguments are shifted
by two. Finally the modulus-product powers use a previous factorial
output, which is at least one. This is an induction through the acyclic
arithmetic-macro dependencies. The supplied audit checks that dependency
induction at all $55$ actual call sites. Theorem
\ref{thm:positive-power11} is not asserted for zero-base or
zero-exponent inputs.

\subsection{The inherited direct source and its bounded compiler}

For clarity, the full semantic chain is recorded here. The exact C222
constructors, eight unconditional axiom families, guarded bundle,
modus-ponens branch, and generalization branch are supplied by
\path{direct/direct_calculus.py}. They are literal compositions of
\[
 I(a,b)=2(a+b)(a+b+1)+4b,
\quad M=I+1,\quad A=I+2,\quad E=I+3.
\]
The axiom and inference formulas were audited against the supplied
\path{Code254.lean} and \path{KCode222.lean}. The imported calibration is
\[
 C_{222}(1)\ \longleftrightarrow\
 \mathrm{ZF}\text{ has a finite refutation}
 \ \longleftrightarrow\
 \mathrm{ZFC}\text{ has a finite refutation}.
\]
Its named results are
\texttt{KEndpoint.c222\_iff\_zf\_refutation} and
\texttt{ConstructibleChoice.zf\_refutation\_iff\_zfc\_refutation}.
The immediate source files are included with their pinned provenance.

The numerical calculus has single-conclusion Horn rules, padded to two
premise slots using the freely derivable truth code $0$. Its target
is code $1$. A finite nonnegative multiset of rule instances proves the
target precisely when its conclusion-minus-premise balance is the
target unit vector, after ignoring coordinate $0$. For soundness,
if the target were underivable, sum the balance over all underivable
formulas. Each rule contributes at most zero, whereas the target
contributes one. Completeness counts nodes in a finite proof tree.

For $t$ rule instances with all labels in $[0,m]$, choose
$M=(m+1)^t$. With $m,t\ge1$, the identity
\[
 \prod_{i<t}(1+v_iM)
    =(1+M)\prod_{i<t}(1+p_iM)(1+q_iM)
\]
implies the corresponding identity of polynomials in a formal variable.
Indeed, if the right product had more than $t$ nonconstant factors,
its value would be at least $M^{t+1}$, whereas the left value is
at most $(1+mM)^t<M^{t+1}$. Thus each side has at most $t$
nonconstant factors. Its coefficient sum is at most $(m+1)^t=M$,
and each individual coefficient is strictly below $M$. Base-$M$
uniqueness then proves the polynomial identity, and unique
factorization gives the exact positive-label balance.

One accumulator realizes this equality with polynomial local steps
\[
 S_0=1+M,\quad S_t=1,\qquad
 S_{i+1}(1+v_iM)=S_i(1+p_iM)(1+q_iM).
\]
Completeness expands open obligations in a proof tree; the accumulator
is the product of the affine weights of those obligations. A beta code
stores the finite sequence. The source has thirteen per-row natural
coordinates: three labels, four schema coordinates and six
sequence/quotient/remainder coordinates. No proof-order or cycle test
is omitted without justification: the single-conclusion balance
argument supplies exactly that justification.

The finite bounded block is compiled using the following shared CRT
principle. Write its nonnegative local polynomial as $P=E-F$, where
$E,F$ have nonnegative coefficients. With coordinate bounds $u_j\le V$,
put
\[
 H=E(t,z,u)+F(t,z,u)+t+V+1,\quad T=H!,\quad
 B_*=\prod_{i<t}(1+(i+1)T).
\]
The whole bounded block is equivalent to finitely many existential
conditions
\[
 B_*\mid\binom{A_j}{u_j+1},\qquad
 B_*\mid1+(C+1)T,\qquad
 B_*\mid E(C,z,A)-F(C,z,A).
\]
For soundness, a prime divisor of $1+(i+1)T$ is greater than $H$.
The binomial condition selects a residue in $[0,u_j]$ for each
coordinate, the index condition selects residue $i$, and the size
bound turns the local congruence into equality. For completeness,
the moduli $1+(i+1)T$ are pairwise coprime. The CRT supplies all
residue classes, and $(u_j+1)!\mid T$ permits cancellation of the
falling-factorial denominator. This proves both directions without
treating factorial or binomial notation as polynomial syntax.

The builder eliminates the product, binomial and factorial operations
by the explicit arithmetic identities and finite graphs retained in
\path{direct/upstream/parent_crt.tex} and
\path{direct/upstream/parent_primitives.tex}. Their only modified
primitive is Theorem~\ref{thm:positive-power11}. For example,
\[
 \binom nk\text{ is digit }k\text{ of }(2^n+2)^n
       \text{ in base }2^n+1,
\]
so three power graphs and four quotient/remainder coordinates suffice.
The factorial identity used is
\[
 n!=\left\lceil
 \frac{(\rho+1)^n}{\binom{\rho+n}{n}}
 \right\rceil,\qquad \rho=(n+2)^{n+2}.
\]
For $n\ge2$, the difference between $n!$ and the fraction is in
$[0,1)$, because
\[
 0\le n!-n!\prod_{i=1}^n\frac{\rho+1}{\rho+i}
 <\frac{n!n(n-1)}{2(\rho+1)}<1.
\]
The cases $n=0,1$ are immediate. Ceiling, division and bounds are
implemented by ordinary natural slack equations. The complete finite
builder emits all their arithmetic gates and residuals.

\subsection{A final sign saving and exact gate elimination}

The local source $P=E-F$ is a product of sums of squares plus further
squares, hence $E\ge F$ on every natural assignment. Its final
divisibility condition therefore has a natural quotient:
\[
 E=F+B_*q,\qquad q\in\mathbb N.
\]
Replacing the two-natural encoding of a signed quotient by this one
coordinate is valid before any local rule is known to hold.

An arithmetic gate $z=x+y$ or $z=xy$ contributes its quadratic
defining residual. A finite conjunction of residual equations is
equivalent to their sum of squares being zero. When a gate is
$z=L$ for an affine polynomial $L$ with nonnegative coefficients in
earlier coordinates, substitution is a bijective elimination of that
natural coordinate. It preserves the quadratic residual degree bound.
The exporter applies this operation only when the expanded affine
support has at most four terms; that threshold controls size and
does not affect correctness. It removes $2\,387$ coordinates and
rows from the retained-affine variant:
\[
 4860-2387=2473,\qquad 4476-2387=2089.
\]
All remaining rows are literally supplied, so this is a completed
ordinary polynomial rather than an unmaterialized compiler proposal.

\subsection{Verification and the logical boundary}

The independent checker imports no builder code. It reconstructs the
entire sum of squares, compares every coefficient with the supplied
expanded polynomial, checks that every declared coordinate occurs,
and verifies the degree, coefficient heights and residual inventories.
It checks all $55$ power blocks for exactly eleven primitive auxiliaries
plus output and exactly six semantic constraint labels. Both final
variants pass in ordinary Python and with optimization enabled.

A separate symbolic audit eliminates every arithmetic gate in one
power block and compares the resulting six polynomial constraints
with an independently written version of
\eqref{eq:op-norms}--\eqref{eq:op-rest}. It also verifies the three
restoration identities, the positive-input induction at every actual
call site, and seven constructive Pell fixtures, including a fixture
with intermediate values exceeding two million bits. Altering the
claimed power output is detected in every fixture. These checks test
the translation and its guards, not the consistency of ZFC.

The new power theorem and all coordinate eliminations have ordinary
proofs. The C222/refutability calibration, the standard Pell-index
theorem and the inherited shared-CRT construction remain identified
mathematical dependencies. Their immediate sources are included; the
full transitive Lean dependency closure was not rebuilt, and no new
end-to-end Lean certificate or object-language ZFC proof term is
claimed. Composing those results with the proved transformations gives
Theorem~\ref{thm:direct2473}; negating refutation existence gives its
consistency formulation.
